%% 第四章

\chapter{测试基准数据构建工具链实现}
\chaptermark{测试基准数据构建工具链实现}

\section{工具链总体流程}

UAV-DualCog 的数据构建并非一次性人工整理，而是由一套多阶段工具链完成。工具链的总体目标是从仿真场景出发，经过场景几何建模、地标资产构建、视频任务生成和图像任务生成，最终导出可评测、可复核、可追踪的结构化测试基准。

\begin{figure}[htbp]
    \centering
    \includegraphics[width=\textwidth]{fig_benchmark_construction_process.pdf}
    \caption{UAV-DualCog四阶段构建流程}
    \label{fig:construction-process}
\end{figure}

如\autoref{fig:construction-process}所示，工具链分为四个阶段：阶段1完成场景扫描与点云融合，阶段2完成地标挖掘、人工复核与语义标注，阶段3完成行为驱动的视频任务构建，阶段4完成证据感知的图像问答任务构建。各阶段与仓库真实实现入口、关键伪代码和代表性逻辑的对应关系见\autoref{app:toolchain-entry}。正文重点说明其数据流、关键约束和任务产物，避免将工程实现直接写成文件清单。

\section{阶段1：场景扫描与点云融合}

阶段1的目标是建立稳定的场景级几何基础。该阶段基于 AerialVLN 模拟器场景和配置文件，在给定边界、采样高度和姿态覆盖范围内进行多位姿扫描，同时采集 RGB 图像、语义分割、激光雷达和无人机位姿信息。随后，系统将这些原始数据统一变换到同一场景坐标系中，融合为带有语义和实例属性的场景级点云。

这一阶段的关键作用在于将原始模拟器场景转换为可复用的几何基底。完成阶段1后，每个场景具备稳定的地图边界、全局坐标系、语义点云、扫描预览图和面积统计信息。后续阶段不需要重复依赖局部图像检测，而可以直接在场景级几何表示上开展地标候选聚合、可见性判断和视角采样。

从工程实现上，阶段1围绕场景配置解析、模拟器桥接采集和场景级产物落盘三个环节组织。系统统一读取场景配置，调用模拟器桥接模块完成多位姿采样，并将点云、日志和缓存结果写入标准化目录结构中。具体实现入口与伪代码见\autoref{app:toolchain-entry}。

\section{阶段2：地标挖掘、复核与语义标注}

阶段2基于阶段1输出的语义与实例点云构建地标资产。系统首先按照实例信息聚合候选对象，形成初始地标候选池；随后根据候选对象三维包围盒和实例中心，计算八个标准方向的拍摄位姿，包括前、右前、右、右后、后、左后、左、左前等视角，并生成对应 RGB 视图和点云视图。

自动候选生成之后，项目引入人工复核机制。复核内容包括候选是否保留为有效地标、主视图是否合理、方向标注是否正确、某些视角是否需要剔除、地标框是否有效以及该地标是否具备后续任务构建价值。通过人工复核，阶段2输出形成有效地标清单，作为图像任务和视频任务的共同资产基础。

在语义标注方面，项目使用受约束的自动标注流程，为每个地标补充粗粒度类别、细粒度子类和简短可区分描述。与自由生成式描述不同，该流程通过系统提示词、用户提示模板和固定 JSON 输出格式约束模型生成，使地标语义信息既可读又可解析。该设计为后续任务提示词、样本复核和实验记录提供了统一语义基础。

阶段2的主要实现入口见\autoref{app:toolchain-entry}。该阶段支持实例收集、复核界面、语义标注和有效实例导出等模式。通过复核界面，标注者可以同时查看实例级信息、RGB 视图、点云视图、方向标签和语义描述，从而提高地标质量控制效率。

\section{阶段3：行为驱动的视频任务构建}

阶段3负责视频任务生成。该阶段基于阶段2得到的有效地标资产和预设行为库，为地标绑定飞行任务，并导出轨迹、视频、行为区间和可见性监督信息。

该阶段的一个重要设计是两层飞行行为体系。底层为原子行为，例如渐进接近、环绕飞行、螺旋环绕、升空观察和面状扫描等；上层为复合行为，例如环绕巡检、螺旋巡检、三角巡检和面状测绘巡检等。复合行为由多个原子行为组合而成，更接近真实巡检和测绘任务中的飞行模式。

阶段3实现包括任务模板参数确定、轨迹生成、碰撞检查、安全距离约束、半径和高度修正、并行调度、关键帧与时间区间监督导出等步骤。最终，该阶段保留高分辨率原始帧和任务级监督文件，同时导出面向测试基准的 1080P H.264 视频和结构化任务清单，使视频任务既适合模型调用，又保留足够的时间定位监督。

从实现结构看，阶段3由轨迹生成与任务组织两个子流程共同承担，具体入口和关键伪代码见\autoref{app:toolchain-entry}。前者负责生成飞行轨迹、视频和过程记录，后者负责组织任务、实验和指标聚合。核心视频任务分别覆盖自我状态认知下的飞行行为识别与时间定位，以及环境情况认知下的地标可见次数与时间定位。

\section{阶段4：证据感知图像问答任务构建}

阶段4负责图像任务生成。该阶段的核心思想不是随机选取截图组成问答，而是围绕地标中心参考视图和任务相关第一视角观测图进行结构化样本生成。自我状态认知和环境情况认知两类图像任务共享部分地标视图资产，但分别围绕“自身位置与未来观察”和“目标方向与动作决策”展开问题设计。

图像任务生成中，系统同时考虑语义答案和空间证据。对于需要定位的任务，样本除问题、选项和答案外，还包含标准化目标框。对于环境情况认知任务，系统会调整观测位姿和相机朝向，直到目标完整进入画面且目标框有效为止。最终输出不仅包括任务清单，也包括重新渲染所需辅助文件，以便在不改变样本语义结构的前提下刷新媒体资源。

阶段4的主要实现入口见\autoref{app:toolchain-entry}。该阶段支持图像 QA 生成、任务渲染、模型调用、结构化输出解析和指标汇总。四类图像任务分别覆盖自我状态认知与环境情况认知下的两个细分方向，并统一使用选项答案和目标框作为输出。

\section{统一管线与配套材料整理}

为了降低多阶段构建带来的复杂度，项目实现了统一管线编排。仓库中的 \texttt{task\_pipeline.py} 支持跨场景的数据筛选、任务生成、渲染、实验和分析流程；\texttt{pipeline\_common.py} 则统一桥接配置、场景目录解析、阶段目录解析、完整性校验和阶段日志。典型场景目录包含 \texttt{pcd\_map}、\texttt{landmarks\_raw}、\texttt{landmarks\_review}、\texttt{landmarks}、\texttt{stage3\_tasks} 和 \texttt{logs} 等子目录。

配套材料方面，项目整理了 README、数据集说明、使用说明、实验榜单、分析结果和配置模板。这些材料分别对应基准介绍、数据使用、模型评测和结果分析等环节，用于支撑测试基准的复核、复现实验和后续扩展。

\section{本章小结}

本章介绍了 UAV-DualCog 测试基准的四阶段构建工具链。阶段1提供场景几何基础，阶段2提供复核后的地标资产，阶段3生成视频任务和时间监督，阶段4生成图像任务和空间证据。通过共享配置、统一路径和结构化输出，该工具链既支撑了当前测试基准的稳定构建，也已经实际支撑 UAV-DualCog-Train 训练集的数据导出。工具链实现入口与关键伪代码见\autoref{app:toolchain-entry}。
